\chapter{Uang dan Pengukuran}\label{ch:g2:measure}

Penggaris, timbangan dapur, teko, uang logam, jam dan kalender: bab
ini berbicara tentang bilangan dalam kehidupan sehari-hari ---
berapa panjangnya, berapa beratnya, berapa isinya, berapa harganya,
dan pukul berapa sekarang.

\section{Panjang, massa, kapasitas}

\begin{definition}[Satuan sehari-hari]\label{def:g2:measure:units}
\begin{itemize}
  \item \emph{Panjang}: sentimeter (cm) untuk benda kecil, meter (m)
        untuk benda besar, dengan $1$ m $= 100$ cm;
  \item \emph{massa}\index{massa}: gram (g) untuk benda ringan, kilogram (kg)
        untuk benda berat, dengan $1$ kg $= 1000$ g;
  \item \emph{kapasitas}\index{kapasitas} (berapa banyak isi sebuah wadah): liter (L).
\end{itemize}
Satuan yang lebih kecil dan aturan pastinya menyusul pada \cref{ch:g3:measure}.
\end{definition}

\begin{example}[Satuan yang mana?]\label{ex:g2:measure:choose}
Sebatang krayon: kira-kira $9$ cm. Sebuah pintu: kira-kira $2$ m.
Sebuah apel: kira-kira $150$ g. Sekantong besar tepung: $1$ kg.
Sebotol susu: $1$ L. Tebaklah satuannya sebelum mengukur --- itu
mencegah jawaban konyol seperti ``pintunya $2$ cm''.
\end{example}

\section{Uang}

\begin{method}[Menyusun sejumlah uang]\label{met:g2:measure:money}
Dengan uang logam dan uang kertas $1$, $2$, $5$, $10$, $20$:
\begin{enumerate}
  \item mulailah dari yang terbesar yang masih muat, lalu isilah
        sisanya;
  \item untuk menyusun $27$: satu $20$, satu $5$, dan satu $2$;
  \item sering ada beberapa cara --- $27$ juga sama dengan
        $10 + 10 + 5 + 2$.
\end{enumerate}
\end{method}

\begin{omfigure}
\begin{tikzpicture}[scale=0.9]
  \draw[thick, fill=omProp!20] (0,0) rectangle (1.3,0.7);
  \node[font=\small] at (0.65,0.35) {$20$};
  \node[font=\large] at (1.9,0.35) {$+$};
  \draw[thick, fill=omMeth!20] (2.85,0.35) circle (0.35);
  \node[font=\small] at (2.85,0.35) {$5$};
  \node[font=\large] at (3.8,0.35) {$+$};
  \draw[thick, fill=omDef!20] (4.7,0.35) circle (0.3);
  \node[font=\small] at (4.7,0.35) {$2$};
  \node[font=\large] at (5.6,0.35) {$=$};
  \node[font=\normalsize] at (6.5,0.35) {$27$};
\end{tikzpicture}

{\small Menyusun $27$ dengan satu uang kertas dan dua uang logam ---
dan masih ada cara lain, seperti $10 + 10 + 5 + 2$.}
\end{omfigure}

\begin{example}[Kembalian]\label{ex:g2:measure:change}
Sebuah buku berharga $14$; kamu membayar dengan $20$. Kembaliannya
$20 - 14 = 6$, ditemukan dengan membilang maju: $14 \to 20$ adalah
$6$ (\cref{met:g2:subtraction:countup}).
\end{example}

\section{Waktu}

\begin{method}[Setengah jam dan seperempat jam]\label{met:g2:measure:clock}
Jarum panjang menunjukkan menit setelah jam bulat:
\begin{enumerate}
  \item lurus ke atas ($12$): tepat pada jamnya --- $3{:}00$;
  \item seperempat putaran ke kanan ($3$): lewat seperempat ---
        $3{:}15$;
  \item lurus ke bawah ($6$): setengah jam lewat --- $3{:}30$, yang
        dibaca ``pukul setengah empat'';
  \item tiga perempat putaran ($9$): seperempat menjelang jam
        berikutnya --- $3{:}45$, ``pukul empat kurang seperempat''.
\end{enumerate}
\end{method}

\begin{omfigure}
\begin{tikzpicture}[scale=0.8]
  \foreach \x/\hh/\mm/\lbl in {0/3/3/{3:15}, 4/3/6/{3:30},
                                8/3/9/{3:45}} {
    \begin{scope}[xshift=\x cm]
    \draw[thick] (0,0) circle (1.2);
    \foreach \h in {1,...,12} {
      \node[font=\tiny] at ({90-\h*30}:0.98) {\h};
    }
    \draw[very thick, omDef] (0,0) -- ({90-(\hh+\mm/12)*30}:0.55);
    \draw[thick, omThm] (0,0) -- ({90-\mm*30}:0.95);
    \node[below, font=\small] at (0,-1.45) {\lbl};
    \end{scope}
  }
\end{tikzpicture}

{\small Lewat seperempat, setengah jam lewat, kurang seperempat:
jarum panjang merah di angka $3$, angka $6$, dan angka $9$.}
\end{omfigure}

\begin{example}[Kalender]\label{ex:g2:measure:calendar}
Satu tahun punya $12$ bulan, berurutan: Januari, Februari, Maret,
April, Mei, Juni, Juli, Agustus, September, Oktober, November, dan
Desember. Kebanyakan bulan punya $30$ atau $31$ hari; Februari punya
$28$ (atau $29$). Sebuah tanggal menyebut hari dan bulannya:
``tanggal $5$ Maret''.
\end{example}

\section{Latihan}

\begin{exercise}[$\star$]\label{exo:g2:measure:1}
Pilihlah satuannya (cm, m, g, kg atau L): sebatang pensil; panjang
ruang kelas; sehelai bulu; sebuah tas sekolah; sebotol air.
\end{exercise}

\begin{exercise}[$\star$]\label{exo:g2:measure:2}
Ubahlah: $2$ m menjadi cm; $300$ cm menjadi m; $1$ kg menjadi g.
\end{exercise}

\begin{exercise}[$\star$]\label{exo:g2:measure:3}
Susunlah \omterm{def:g1:addition:def}{jumlah} uang ini dengan logam dan uang kertas: $16$; \quad
$23$; \quad $35$. Berikan satu cara untuk masing-masing.
\end{exercise}

\begin{exercise}[$\star$]\label{exo:g2:measure:4}
Hitunglah kembaliannya: bayar $10$ untuk mainan seharga $6$; bayar
$20$ untuk buku seharga $17$; bayar $50$ untuk permainan seharga $32$.
\end{exercise}

\begin{exercise}[$\star$]\label{exo:g2:measure:5}
Bacalah waktunya: jarum panjang di $6$, jarum pendek antara $4$ dan
$5$; jarum panjang di $3$, jarum pendek tepat setelah $8$.
\end{exercise}

\begin{exercise}[$\star$]\label{exo:g2:measure:6}
Gambarlah jam yang menunjukkan pukul setengah tiga. Ke mana setiap jarum menunjuk?
\end{exercise}

\begin{exercise}[$\star$]\label{exo:g2:measure:7}
Sebutkan bulannya: sesudah Mei; sebelum Januari; bulan hari ulang
tahunmu.
\end{exercise}

\begin{exercise}[$\star$]\label{exo:g2:measure:8}
Lia membeli minuman seharga $3$ dan kue seharga $4$. Ia membayar
dengan uang kertas $10$. Berapa kembaliannya?
\end{exercise}

\begin{exercise}[$\star$]\label{exo:g2:measure:9}
Mana yang lebih berat: sekantong gula $1$ kg atau $800$ g apel?
Berapa gram \omterm{ex:g1:subtraction:difference}{selisihnya}?
\end{exercise}

\begin{exercise}[$\star\star$]\label{exo:g2:measure:10}
Susunlah $40$ dengan tepat empat uang logam atau uang kertas
(dipilih dari $1$, $2$, $5$, $10$, $20$). Carilah dua cara berbeda.

\end{exercise}
